% !TEX root = ../main.tex
%-----------------------------------------------------------------------
%  The recommendation on kappa_v as a network input, its consequences for
%  training data, and the next steps.
%  \input by ../main.tex -- not compiled on its own.
%-----------------------------------------------------------------------
\hd{Recommendation: include \kv{} as an input to the neural network.}
The surrogate design already reserved a place for \kv{} among its input
fields, pending this experiment's verdict. The verdict is in, and the
recommendation is unambiguous, for four reasons.

\begin{enumerate}[nosep,leftmargin=1.6em]
  \item \textbf{The sensitivities measurably change within the plausible
        mixing range.} A factor-of-two change in \kv{} --- well within the
        uncertainty ocean models carry for this parameter --- shifts
        sensitivity amplitudes by 20--40\,\% and reshapes the
        mixing-sensitivity field itself (Finding~2). A network without \kv{}
        as an input cannot represent this variation at all: it would collapse
        distinct answers onto one.
  \item \textbf{The dependence is nonlinear from the smallest step tested.}
        Because linear extrapolation in \kv{} already fails at a factor of two
        (Finding~4), the alternative of training at the reference and
        correcting predictions afterwards with the adjoint's own gradient is
        not viable. The nonlinearity must be learned, which requires \kv{}
        in the inputs and training data that spans its range.
  \item \textbf{The output most worth predicting is the most
        \kv-dependent.} The gradient $\partial J/\partial\kv$ --- the field a
        mixing-aware surrogate exists to deliver --- varies fastest and
        changes sign across the ensemble (Finding~2). Predicting it from
        state alone, without knowing the mixing, is not a well-posed task.
  \item \textbf{The cost of inclusion is negligible; the cost of omission is
        structural.} \kv{} enters as one more input channel (it is a
        three-dimensional field in general, even though this ensemble varied
        it uniformly). Adding it later, after an architecture and training
        corpus were built without it, would mean regenerating both.
\end{enumerate}

\hd{What the results change about generating training data.}
The unexpected constraint is numerical, not physical: adjoint stability, not
computer time, now bounds the training set. Full-length (five-year)
sensitivity maps exist only for $0.5\times$--$2\times$ (plus the flagged
$16\times$); the unstable members contribute only their early, still-accurate
portions --- 62\,\% of the ensemble's 5-day records survive in total. Three
routes forward, not mutually exclusive: (i)~rerun unstable members with the
stabilised (\code{adjViscBoost}) adjoint, first measuring the stabiliser's
own bias on a member where the plain answer is valid; (ii)~use shorter
adjoint windows, which stay inside every member's stable span;
(iii)~train on lead-limited records, using each member's measured departure
time as its label horizon. Two further practical points: the model's internal
variability (2--4\,\% of $J$) is the natural floor for surrogate accuracy ---
pushing integrated errors below it is fitting noise --- and six of the
fourteen control gradients (initial velocities and the four wind controls)
are currently written as zeros and must be wired or dropped before anything
trains on them.

\hd{Next steps, in priority order.}
(1)~A $1\times$ control member, to isolate the adjustment effect of
Finding~1 and double the noise-floor sample. (2)~The stabilised rerun of the
worst unstable member ($4\times$), paired with a stabilised rerun of the
healthy $2\times$ member as the bias control. (3)~Two members at $0.7\times$
and $1.4\times$, inside the proven-stable range, to resolve the curvature
that Finding~4 exposed. (4)~The long-planned repair of the model's built-in
gradient check at a point of large sensitivity --- more urgent now that
Finding~4 shows large perturbations cannot validate the gradient. (5)~Resolve
the six zero-valued control gradients.

\hd{Where the full analysis lives.}
Every number here is computed in a suite of seven documented notebooks at
\code{analyses/DINO\_1deg/03\_adjoint/05\_kappa\_v\_ensemble/} in the project
repository (figures, summary tables, and the animations listed above sit
beside it on cluster storage), and is restated with sources in the master
plan's Part~I Results section, which remains the authoritative record.
